\documentclass[11pt]{article}
\usepackage[a4paper,margin=27mm]{geometry}
\usepackage[T1]{fontenc}
\usepackage{lmodern,amsmath,amssymb,amsthm,booktabs,array,longtable}
\usepackage[colorlinks=true,linkcolor=blue,citecolor=blue,urlcolor=blue]{hyperref}
\setlength{\emergencystretch}{3em}
\newtheorem{theorem}{Theorem}
\newtheorem{lemma}[theorem]{Lemma}
\newtheorem{corollary}[theorem]{Corollary}
\newcommand{\N}{\mathbb N}
\newcommand{\Z}{\mathbb Z}
\newcommand{\coef}[2]{[z^{#1}]#2}
\title{Modular recurrences and alternative proofs of\\
Blitvi\'c, Kammoun and Steingr\'imsson's\\$14325$ packing congruences}
\author{Haobo Ma\\\small ChronoAI Pte Ltd, Singapore
\and Wenlin Zhang\\\small National University of Singapore\\
\small ChronoAI Pte Ltd, Singapore}
\date{19 November 2026}
\begin{document}
\maketitle
\begin{abstract}
Starting from Blitvi\'c, Kammoun and Steingr\'imsson's Jacobi-fraction
theorem (Part~II, arXiv:2611.12829v1, Theorem~9.2, 9 November 2026),
we extend their congruence observations in Part~I (arXiv:2610.29846v1,
Section~7, item~11, p.~49) to finite linear recurrences modulo every
integer coprime to six.
In particular, the prescribed $14325$ packing count is one modulo
thirty-five for an odd number of windows and six for an even number.
An alternative proof of the original modulo-ten and four-divisibility
laws uses Wallner's rank recurrence and Catalan numbers.
For prescribed mixtures of overlaps one and two we give product
congruences and characterize exactly when the comparison poset has
a separator.
\end{abstract}

\section{The enumeration and the congruences}
For $j\geq1$, let $A_j$ count the permutations $w$ of
$\{1,\ldots,3j+2\}$ for which each window
\[
 (w_{3t+1},w_{3t+2},w_{3t+3},w_{3t+4},w_{3t+5}),
 \qquad 0\leq t<j,
\]
has relative order $14325$. These are precisely the $(2,j)$-packings
of Definition~2.6 of Blitvi\'c, Kammoun and Steingr\'imsson~\cite{BKS}.
Additional occurrences are allowed, as in their Remark~2.7.
Their Section~7, item~11 records two numerical observations about $A_j$,
the first through at least $j=23$. Their later Jacobi-fraction
theorem~\cite[Theorem~9.2]{BKSII} already implies both congruence
laws. We derive from it a finite modular model for every modulus
coprime to six (Theorem~\ref{thm:modular}), including explicit laws
modulo thirty-five and one hundred forty. We also give a self-contained
proof of the original laws from the recurrence recorded by Wallner
in 2024~\cite{OEIS}, and determine the separators for mixed overlaps.

Write $C_j=\binom{2j}{j}/(j+1)$ for the Catalan numbers, and let $s_2(n)$
be the number of ones in the binary expansion of $n$.
\begin{theorem}\label{thm:main}
We have $A_1=1$. For every $j\geq1$, $A_j\equiv1\pmod5$; for $j\geq2$,
\begin{equation}\label{eq:cat}
 A_j\equiv3C_j\pmod4.
\end{equation}
Consequently, for every $j\geq1$,
\begin{align}
 A_j&\equiv
 \begin{cases}1&j+1=2^k\text{ for some }k\geq1,\\
 6&\text{otherwise}\end{cases}\pmod {10},\label{eq:ten}\\
 4\nmid A_j&\quad\Longleftrightarrow\quad
 j+1=2^i+2^p\text{ for some }0\leq p\leq i.\label{eq:four}
\end{align}
Equal exponents and exponent zero are included in~\eqref{eq:four}.
\end{theorem}

The underlying grid sequence was already enumerated by Pan~\cite{Pan},
page~3, in 2016 and is OEIS A274644~\cite{OEIS}. Pan's recursive algorithm on page~2
and the packing-poset framework in Sections~2.3 and~4.2 of~\cite{BKS}
give the enumerative setting. Wallner's formula in that entry, dated
13 February 2024, is exactly~\eqref{eq:transfer}--\eqref{eq:A} after
an index shift. We prove its rank-insertion interpretation before
analyzing the two moduli. The Catalan digit-sum criterion
is classical: Deutsch and Sagan~\cite[Theorem~2.1]{DS} give a
combinatorial proof. We include an arithmetic proof of that auxiliary
criterion below; the assertions about $A_j$ require the grid congruence.

\section{A rank-insertion recurrence}
The insertion below is also the reversible insertion in
\cite[Lemma~9.3, equation~(9.11)]{BKSII}: its state $q$ corresponds
to $k=q+2$ here. We count continuations from a state, rather than
extensions arriving in that state.
Use zero-based positions, and put $a_t=w_{3t}$, $b_t=w_{3t+1}$,
$c_t=w_{3t+2}$. The window condition is exactly
\begin{equation}\label{eq:chain}
 a_t<a_{t+1}<c_t<b_t<b_{t+1}\qquad(0\leq t<j).
\end{equation}
Both $a_0$ and $b_j$ are comparable with every entry: the $a_t$ increase,
the $b_t$ increase, and every $c_t$ lies between the adjacent entries
in~\eqref{eq:chain}. Thus $a_0$ is the least entry and $b_j$ the greatest.
Delete these two forced entries and standardize to ranks $0,\ldots,3j-1$.
The remaining entries, indexed by
\[
 x_t=a_{t+1},\qquad y_t=c_t,\qquad z_t=b_t\qquad(0\leq t<j),
\]
satisfy precisely
\begin{equation}\label{eq:grid}
 x_t<y_t<z_t,\qquad x_t<x_{t+1},\qquad z_t<z_{t+1}.
\end{equation}
Conversely, these inequalities, with a new least $a_0$ and greatest
$b_j$, recover every comparison in~\eqref{eq:chain}, including the
first and last windows. Hence $A_j$ is the number of rank assignments
to~\eqref{eq:grid}; there is no multiplicative factor.

For a grid of $j$ columns let $N=3j$, let $h$ be the rank of $x_{j-1}$,
and define its state to be $k=N-h$. Append a new column by inserting
$z_j$ as the greatest rank, then inserting $y_j$ at rank $b$ and $x_j$
at rank $a$, where
\begin{equation}\label{eq:insert}
 h<a\leq b\leq N.
\end{equation}
Here inserting a rank shifts every existing rank at least that large
up by one. The new ranks of $x_j,y_j,z_j$ are $a,b+1,N+2$.
These inequalities give exactly~\eqref{eq:grid} for the new column.
Conversely $z_j$ must be the greatest entry in any grid, since all
columns lie below their own $z_t$ and the $z_t$ increase.
Deleting the last column and standardizing therefore reverses this
construction uniquely.

Set $r=N-a+1$. Then $1\leq r\leq k$, the number of choices for $b$
is $r$, and the new state is $N+3-a=r+2$.
For $n,k\geq0$ define
\begin{equation}\label{eq:transfer}
 F_0(k)=1,\qquad F_{n+1}(k)=\sum_{r=1}^{k}rF_n(r+2).
\end{equation}
For every actual state $k$, induction on $n$ shows that $F_n(k)$ is
the number of extensions by $n$ columns. There is one single-column
grid, with state three. Consequently
\begin{equation}\label{eq:A}
 A_j=F_{j-1}(3)\qquad(j\geq1).
\end{equation}
Values of $F_n$ at auxiliary arguments below three are fixed
algebraically by~\eqref{eq:transfer}.

\section{Reduction modulo five}
Let $Tf(k)=\sum_{r=1}^k rf(r+2)$, and regard functions as taking values
in $\Z/5\Z$. Define the period-five functions $P,Q$ by
\[
\begin{array}{c|rrrrr}
 k\bmod5&0&1&2&3&4\\\hline
 P(k)&0&1&3&1&0\\
 Q(k)&0&1&1&1&0
\end{array}
\]
Then $T1=P$, $TP=Q$, and $TQ=Q$. To check these identities for all
$k$, use $Tf(0)=0$ and
$Tf(k+1)-Tf(k)=(k+1)f(k+3)$.
The differences of $P$ are $(1,2,3,4,0)$; those of $Q$ are
$(1,0,0,4,0)$, which agree with both $(k+1)P(k+3)$ and
$(k+1)Q(k+3)$. Induction on $k$ proves the three identities.
Thus $F_0(3)=1$, $F_1(3)=P(3)=1$, and $F_n(3)=Q(3)=1$ for $n\geq2$,
proving $A_j\equiv1\pmod5$.

\section{Reduction modulo four}
All equalities in this section are in $R=\Z/4\Z$ or $R[[z]]$.
First, for all $n,q\geq0$,
\begin{equation}\label{eq:neighbor}
 2F_n(2q+2)=2F_n(2q+1).
\end{equation}
For $n=0$ this is immediate; for $n\geq1$ subtract the two sums
in~\eqref{eq:transfer} to obtain
$(2q+2)F_{n-1}(2q+4)$, whose double vanishes in $R$.
In particular,
\[
 (2q+2)F_n(2q+4)=-(2q+2)F_n(2q+3).
\]
Writing $G_n(k)=F_n(2k+1)$, this pairs the even-weight terms in the
transfer sum with the preceding odd-weight terms, giving
\begin{equation}\label{eq:odd}
 G_{n+1}(k)=(2k+1)G_n(k+1)-\sum_{r=1}^kG_n(r).
\end{equation}
Indeed each pair with weights $2r-1,2r$ contributes
$-F_n(2r+1)=-G_n(r)$; the unpaired last weight is $2k+1$.

Let $C=\sum_{j\geq0}C_jz^j$, so $C=1+zC^2$, and set
\[
 V=3C^2-2,\qquad U=1+zV.
\]
Squaring $C=1+zC^2$ gives $C^2=1+2zC^2+z^2C^4$.
Substitution, with $4=0$ in $R$, gives
\begin{equation}\label{eq:V}
 V+2zV+z^2V^2=1.
\end{equation}
For completeness, the left side minus one simplifies to
$3C^2+2zC^2+z^2C^4+1$, which is zero by the preceding identity.
Therefore
\begin{equation}\label{eq:UV}
 U=(1+2z)V+zUV,\qquad V=UV-zV^2.
\end{equation}
Define
\[
 H_{2m}=UV^m,\qquad H_{2m+1}=V^{m+1}\qquad(m\geq0).
\]
Their constant coefficients are all one, $H_0=1+zH_1$, and
\begin{equation}\label{eq:adjacent}
 H_k=(1+(2k+2)z)H_{k+1}-(2k+3)zH_{k+2}.
\end{equation}
For even $k$, this is the first identity in~\eqref{eq:UV} multiplied
by $V^{k/2}$; for odd $k$, it is the second multiplied by
$V^{(k-1)/2}$. Beginning with $k=0$ and using~\eqref{eq:adjacent}
inductively yields
\begin{equation}\label{eq:H}
 H_k=1+z\left((2k+1)H_{k+1}-\sum_{r=1}^kH_r\right).
\end{equation}
Explicitly, rearrange~\eqref{eq:adjacent} to solve for
$(1+(2k+2)z)H_{k+1}$, and substitute~\eqref{eq:H} for $H_k$;
the resulting identity is~\eqref{eq:H} with $k+1$ in place of $k$.
Comparison of coefficients, inducting on degree simultaneously for
all $k$, now identifies
\[
 H_k=\sum_{n\geq0}G_n(k)z^n.
\]
In particular $H_1=V=3C^2-2$, so for $j\geq2$,
\[
 A_j=\coef{j-1}{V}=3\coef{j-1}{C^2}=3C_j\quad\text{in }R.
\]
This proves~\eqref{eq:cat}. The constant coefficient of $V$ is one,
which accounts separately for $A_1=1$.

\section{Binary exceptions and the prior Jacobi fraction}
Let $v_2(m)$ denote the exponent of two in a positive integer $m$.
Counting multiples of each power of two in $m!$ gives
$v_2(m!)=\sum_{r\geq1}\lfloor m/2^r\rfloor=m-s_2(m)$.
Hence
\[
 v_2(C_j)=s_2(j)-v_2(j+1)=s_2(j+1)-1.
\]
The last identity follows because adding one changes the trailing
$v_2(j+1)$ binary ones of $j$ to zeros and introduces one new one.
Thus $C_j$ is odd exactly when $j+1$ is a power of two.
Together with $A_j\equiv1\pmod5$ and~\eqref{eq:cat}, this proves
\eqref{eq:ten}, including $j=1$ directly.
Similarly, $4\nmid C_j$ exactly when $s_2(j+1)\leq2$.
For $j+1\geq2$, an integer with two binary ones is a sum of two
distinct powers of two, while one with one binary one is
$2^i+2^i$ for some $i\geq0$. The converses are immediate, proving
\eqref{eq:four}; $j=1$ corresponds to $i=p=0$.

The exact Catalan residue, due to Eu, Liu and Yeh~\cite[Theorem~2.3]{ELY},
is $1,2,0$ modulo four according as $s_2(j+1)$ is one, two, or
at least three. The valuation proves the last two cases. For the
first, write $j=2^k-1$. Pairing terms in the Catalan recurrence gives
\[
 C_{2q+1}=C_q^2+2\sum_{i=0}^{q-1}C_iC_{2q-i},\qquad
 q=2^{k-1}-1.
\]
A product in this sum could be odd only if
$i+1=2^a$ and $2q-i+1=2^b$, forcing
$2^a+2^b=2^k$ and hence $i=q$. Thus every paired term is even,
and $C_j\equiv C_q^2\equiv1\pmod4$. The case $k=1$ is immediate.
There is no Catalan residue three. Consequently, for $j\geq2$,
\[
 A_j\equiv\begin{cases}
 11&s_2(j+1)=1,\\
 6&s_2(j+1)=2,\\
 16&s_2(j+1)\geq3
 \end{cases}\pmod{20}.
\]
This uses the Chinese remainder theorem for moduli four and five;
$A_1=1$ is the separate initial term.

\paragraph{Consequences of the prior enumeration.}
Theorem~9.2 of~\cite{BKSII}, with the packing interpretation
in its Lemma~9.3, gives
\[
 \mathcal A(z)=\sum_{j\geq1}A_jz^{j-1}
 =\cfrac{1}{1-6z-\cfrac{35z^2}{1-18z-\cfrac{143z^2}{1-30z-\ddots}}},
\]
with Jacobi coefficients $\alpha_n=12n+6$ for $n\geq0$ and
$\beta_n=36n^2-1$ for $n\geq1$. Modulo five, $\beta_1=0$ and $\alpha_0=1$,
so $\mathcal A=(1-z)^{-1}$. Modulo four every $\alpha_n=2$
and every $\beta_n=3$, so all tails coincide and
$\mathcal A=1+2z\mathcal A+3z^2\mathcal A^2$.
This equation has a unique constant-one solution, by induction on
degree. Equation~\eqref{eq:V} identifies it as $3C^2-2$.
Thus the prior Jacobi fraction gives another proof of both
congruences in Theorem~\ref{thm:main}. Reduction of the formal
continued fraction is valid coefficient by coefficient: its finite
truncations stabilize in each degree and all denominators have
constant term one.

\section{Finite modular models}
The vanishing of a Jacobi down-step gives more than the modulo-five
law. Set $a_n=A_{n+1}$ for $n\geq0$.
\begin{theorem}\label{thm:modular}
Let $m\geq2$ with $\gcd(m,6)=1$, and let
\[
 r=\min\{s\geq1:m\mid 36s^2-1\}.
\]
This integer exists, with $r\leq(m-1)/6$ when $m\equiv1\pmod6$ and
$r\leq(m+1)/6$ when $m\equiv5\pmod6$.
Let $J_r$ be the $r$ by $r$ matrix, indexed from zero, with
\[
 (J_r)_{h,h}=12h+6,\quad (J_r)_{h,h+1}=1,\quad
 (J_r)_{h,h-1}=36h^2-1,
\]
where entries outside the matrix are omitted and all other entries
are zero. Over $\Z/m\Z$,
\begin{equation}\label{eq:finite}
 \sum_{n\geq0}a_nz^n
 =e_0^{\mathsf T}(I-zJ_r)^{-1}e_0
 =\frac{\det\bigl((I-zJ_r)_{\widehat0,\widehat0}\bigr)}{\det(I-zJ_r)}.
\end{equation}
The numerator deletes row and column zero; the empty determinant is one. If
$D(z)=\det(I-zJ_r)=1+d_1z+\cdots+d_rz^r$, then
\begin{equation}\label{eq:modrec}
 a_n+d_1a_{n-1}+\cdots+d_ra_{n-r}=0\pmod m
 \qquad(n\geq r).
\end{equation}
\end{theorem}
\begin{proof}
The indicated choices $s=(m-1)/6$ or $(m+1)/6$ satisfy
$6s\equiv-1$ or $1\pmod m$, proving existence and the bound.
The Jacobi fraction counts weighted Motzkin excursions: an up-step
has weight one, a level-step at height $h$ has weight $12h+6$, and
a down-step from height $h$ has weight $36h^2-1$.
For completeness, the generating function for excursions above height
$h$ is $(1-(12h+6)z-(36(h+1)^2-1)z^2T_{h+1})^{-1}$, by the
first-return decomposition. This is the Jacobi fraction term by term.
Any excursion reaching height $r$ must return through a down-step
of weight $36r^2-1$, so its contribution is zero modulo $m$.
The remaining excursions stay in $0,\ldots,r-1$ and are counted by
$(J_r^n)_{0,0}$. Their geometric matrix series and the adjugate formula
give~\eqref{eq:finite}. The denominator has constant term one and
degree at most $r$, and its cofactor has degree at most $r-1$.
Multiplying the series by $D$ proves~\eqref{eq:modrec}, including
the possible vanishing of its last coefficient.
\end{proof}

\begin{corollary}\label{cor:35}
For all $j\geq1$,
\[
 A_j\equiv6^{j-1}\equiv
 \begin{cases}1&j\text{ odd},\\6&j\text{ even}\end{cases}\pmod{35}.
\]
For $j\geq2$, this strengthens the modulo-twenty consequence to
\begin{equation}\label{eq:140}
 A_j\equiv36\cdot6^{j-1}+35C_j\pmod{140}.
\end{equation}
\end{corollary}
\begin{proof}
For $m=35$, the first down-weight is $\beta_1=35=0$, so $r=1$ and
$\sum a_nz^n=(1-6z)^{-1}$. Since $6^2\equiv1\pmod{35}$, the
parity formula follows. The right side of~\eqref{eq:140} is
$6^{j-1}$ modulo $35$ and $3C_j$ modulo four; apply the Chinese
remainder theorem and~\eqref{eq:cat}.
\end{proof}

For example, $r=4$ for $m=25$, and the finite model is
\begin{equation}\label{eq:25}
 \sum_{n\geq0}a_nz^n=
 \frac{1+10z+15z^2+15z^3}{1+4z+20z^2+20z^3+15z^4}\pmod{25}.
\end{equation}
This follows by expanding the two determinants in~\eqref{eq:finite};
the successive down-weights are $35,143,323,575$, the last divisible
by $25$. Thus the displayed coefficients give a four-term modular
recurrence, rather than a finite-prefix extrapolation.

\section{Mixed overlaps and separators}
\begin{corollary}\label{cor:mixed}
For a nonempty sequence of prescribed $14325$ windows with successive
overlaps one or two, let $j_1,\ldots,j_a$ be the numbers of windows
in its maximal blocks joined by overlaps two. Its count $B$ satisfies
\[
 B=\prod_{\ell=1}^a A_{j_\ell},\qquad
 B\equiv6^{\sum_\ell(j_\ell-1)}\pmod{35}.
\]
Moreover $B\equiv1\pmod5$, and
\[
 B\equiv\begin{cases}
 1&\text{every }j_\ell+1\text{ is a power of two},\\
 6&\text{otherwise}
 \end{cases}\pmod {10},\qquad
 4\nmid B\ \Longleftrightarrow\
 \sum_{\ell=1}^a\bigl(s_2(j_\ell+1)-1\bigr)\leq1.
\]
\end{corollary}
\begin{proof}
The product formula was proved in \cite[Proposition~5, p.~9]{BKS23}
and restated in \cite[Proposition~6.1]{BKS}. Partition a sequence of prescribed
$14325$ windows into its maximal blocks joined by overlaps of size two;
let their numbers of windows be $j_1,\ldots,j_a$, all positive.
Successive blocks overlap in one position. That shared position is
the greatest entry in the left block and the least in the right block,
by~\eqref{eq:chain}. It separates every entry on its left from every
entry on its right. Thus the sets of ranks in the blocks are forced,
and their internal orders can be chosen independently. The count is
$B=\prod_{\ell=1}^a A_{j_\ell}$, the specialization of that proposition.
The residues follow from Theorem~\ref{thm:main} and
Corollary~\ref{cor:35}. For the final equivalence,~\eqref{eq:cat} implies
$\min(2,v_2(A_j))=\min(2,s_2(j+1)-1)$ for $j\geq2$,
and this also holds for $j=1$. Add valuations over the product.
\end{proof}

In Part~I, Section~7, item~5, a separator is a vertex comparable with
every vertex whose deletion leaves a lower and an upper subposet,
neither a chain. The following gives an exact answer within the
mixed-overlap family; it is not a characterization of all patterns.
\begin{corollary}\label{cor:separator}
A prescribed $14325$ packing poset with overlaps one or two has a
separator if and only if at least two of its maximal overlap-two
blocks contain at least two windows. In particular, a pure
overlap-two packing has no separator, for every $j\geq1$.
\end{corollary}
\begin{proof}
In the grid~\eqref{eq:grid}, for $j\geq2$, the only vertices
comparable with all others are its least vertex $x_0$ and greatest
vertex $z_{j-1}$. Indeed $x_i$ is incomparable with $y_{i-1}$ for
$i\geq1$, $y_i$ with $x_{i+1}$ for $i<j-1$, $y_{j-1}$ with
$z_{j-2}$, and $z_i$ with $y_{i+1}$ for $i<j-1$.
These incomparabilities follow directly from the directed relations:
no relation enters an $x$ from a $y$ or $z$, no relation leaves a $z$
for a $y$, and the $x$ and $z$ chains only increase their indices.
Restoring the removed least and greatest entries adds only an
initial and a terminal chain. Thus no vertex splits a single block
into two nonchains. A one-window block is itself a chain.
The preceding product proof expresses a mixed packing as these
blocks joined in order at their extreme vertices. If two blocks
are nonchains, a shared endpoint between them is comparable with
all vertices and has a nonchain on either side, hence is a separator.
If there is at most one nonchain block, a vertex in a joining chain
leaves at least one side a chain. A universal vertex within the
unique nonchain block lies on its initial or terminal chain, so
again at least one side is a chain. This proves both directions.
\end{proof}

\begin{corollary}\label{cor:catalan}
For each pattern--overlap pair
\[
 (1324,2),\ (14253,3),\ (136245,3),\ (142536,4),
\]
\[
 (1376245,3),\ (369247815,5),\ (162738495,7),
\]
the packing count has two-adic valuation $s_2(j+1)-1$ for $j\geq1$.
Thus it is odd exactly when $j+1$ is a power of two, and fails to be
divisible by four exactly when $j+1$ is a sum of two powers of two,
allowing equal exponents.
\end{corollary}
\begin{proof}
Part~I, Section~7, item~6 identifies these Catalan classes by posets
that differ by initial and terminal chains. Their counts are $C_j$,
so the displayed Catalan valuation proves the assertion. This is a
consequence of the cited enumeration, not an additional Wilf equivalence.
\end{proof}

\begin{corollary}\label{cor:equivalent}
For every $j\geq1$, the overlap-two counts for the four patterns
$14325$, $52341$, $25314$ and $41352$ are equal. Each therefore
satisfies all the single-block congruences and finite modular
recurrences above.
\end{corollary}
\begin{proof}
The equivalence of $14325$ and $25314$ is recorded in
\cite[Table~1, p.~11]{BKS23} and \cite[Section~3]{BKS}.
An explicit bijection follows from their inequalities.
For $25314$ they are
\[
 a_{t+1}<a_t<c_t<b_{t+1}<b_t\qquad(0\leq t<j).
\]
Relabel positions by $a'_s=a_{j-s}$, $b'_s=b_{j-s}$ and
$c'_s=c_{j-1-s}$. The inequalities become
$a'_s<a'_{s+1}<c'_s<b'_s<b'_{s+1}$, exactly~\eqref{eq:chain}.
This position relabelling is invertible. Complementing every
permutation value exchanges $14325$ with $52341$ and $25314$ with
$41352$, preserving prescribed starts. The statement concerns
overlap two; it does not assert the mixed-overlap product for $25314$.
\end{proof}

\section{Source coverage and verification}

{\small
\begin{longtable}{p{.40\linewidth}p{.50\linewidth}}
\toprule
Source item (Part I unless indicated)&Coverage\\\midrule
\endhead
Section 7, item 11, first observation&Reproved; prior proof implied by Blitvi\'c, Kammoun and Steingr\'imsson (2026), Theorem 9.2. Equation~\eqref{eq:ten} holds for all $j\geq1$.\\
Section 7, item 11, second observation&Reproved; same prior theorem and Deutsch--Sagan (2006). Equation~\eqref{eq:four} includes zero and equal exponents.\\
Definition 2.6 and Remark 2.7&The enumeration agrees exactly, with additional occurrences allowed.\\
Proposition 6.1; Section 7, item 2&Prior product formula reproved for $14325$; Corollary~\ref{cor:mixed} adds modular consequences. The general varying-overlap question is not covered.\\
Section 7, item 5 (separators)&Not posed for this pattern; obstruction proved: no separator for $14325$ at overlap two. Corollary~\ref{cor:separator} characterizes mixed overlaps; the general classification is not covered.\\
Section 7, item 3 (gap enumeration)&The known grid specialization is~\eqref{eq:transfer}; the search for other classes is not covered.\\
Section 7, item 6 (Catalan classes)&Prior Catalan enumeration cited there; Corollary~\ref{cor:catalan} gives the resulting valuations. The general operations question and shifted-factorial classes are not covered.\\
Section 7, items 1, 4, 7--10&Not covered: general poset restrictions, maxima, Wilf classification, cutoff bounds and other product sequences.\\
OEIS A274644; cross-references A000108 and A274763&Wallner's 2024 recurrence is reproved; A000108 supplies the Catalan comparison. A274763 is a different grid family, not covered.\\
Part II, Theorem 9.2&Prior proof cited and used. Theorem~\ref{thm:modular} and Corollary~\ref{cor:35} are deductions, not a new enumeration.\\
Part I, Table 1 and Section 3; the 2023 Table 1&Prior equivalence reproved in Corollary~\ref{cor:equivalent}; the single-block modular consequences apply to all four patterns.\\
Moduli coprime to six; modulo 140&Not posed in item 11; proved in Theorem~\ref{thm:modular} and~\eqref{eq:140}.\\
\bottomrule
\end{longtable}
}

\paragraph{Verification scope.}
The argument above is a written proof for every $j\geq1$.
The accompanying \texttt{verify.py} independently counts the literal
permutations for $j=1,2$ and their comparison-poset linear extensions
for $1\leq j\leq8$, checks the exact transfer counts and congruences
through $j=256$, and checks the power-series identities through
degree~64. It also checks all prescribed overlap words of sizes one
and two with at most four windows and the Jacobi fraction modulo
twenty through degree~64. Further exact checks compare the modular
matrices for every $2\leq m\leq100$ coprime to six through $j=256$,
the modulo-thirty-five law for $1\leq j\leq256$ and the
modulo-one-hundred-forty law for $2\leq j\leq256$,
\eqref{eq:25} through degree~255, and the separator criterion
for all overlap words with at most seven windows.
The seven listed Catalan classes and the four equivalent patterns
are also checked by comparison-poset enumeration through eight windows.
These finite checks supplement the universal proofs.
In the pinned Lean files~\cite{Core,Transfer}, the declarations
\texttt{count\_mod\_ten} and \texttt{count\_mod\_four\_iff} state
\eqref{eq:ten} and~\eqref{eq:four} for the literal packing predicate;
\texttt{packingGridEquiv} and \texttt{grid\_card\_eq\_row} give the
enumerative correspondence. The declaration
\texttt{row\_catalan\_mod\_four} in~\cite{Transfer} states the
positive-index transfer-row congruence used in~\eqref{eq:cat}.
In that same file, \texttt{row\_mod\_five} states the all-index
modulo-five transfer congruence. The finite modular model,
modulo-thirty-five and modulo-one-hundred-forty laws, mixed-overlap
formulas and separator criterion are written proofs.

\begin{thebibliography}{6}
\bibitem{BKS}
N.~Blitvi\'c, M.~S.~Kammoun and E.~Steingr\'imsson,
\emph{Permutations, Pattern Covers, Posets and Positivity:
Part I: Patterns and Posets}, arXiv:2610.29846v1 (2026),
Definition~2.6, Remark~2.7 and Section~7, item~11, p.~49.
\href{https://arxiv.org/abs/2610.29846v1}{arXiv:2610.29846v1}.
\bibitem{BKSII}
N.~Blitvi\'c, M.~S.~Kammoun and E.~Steingr\'imsson,
\emph{Permutations, Pattern Covers, Posets and Positivity, Part II:
Regular Pattern Packings and Moment Sequences},
arXiv:2611.12829v1 (2026), Theorem~9.2 and Lemma~9.3.
\href{https://arxiv.org/abs/2611.12829v1}{arXiv:2611.12829v1}.
\bibitem{DS}
E.~Deutsch and B.~E.~Sagan,
Congruences for Catalan and Motzkin numbers and related sequences,
\emph{J. Number Theory} \textbf{117} (2006), 191--215.
\href{https://doi.org/10.1016/j.jnt.2005.06.005}{doi:10.1016/j.jnt.2005.06.005};
\href{https://arxiv.org/abs/math/0407326}{arXiv:math/0407326}.
\bibitem{ELY}
S.-P.~Eu, S.-C.~Liu and Y.-N.~Yeh,
Catalan and Motzkin numbers modulo 4 and 8,
\emph{European J. Combin.} \textbf{29} (2008), 1449--1466;
available online 20 August 2007.
\href{https://doi.org/10.1016/j.ejc.2007.06.019}{doi:10.1016/j.ejc.2007.06.019}.
\bibitem{BKS23}
N.~Blitvi\'c, M.~S.~Kammoun and E.~Steingr\'imsson,
A new perspective on positivity in (consecutive) permutation patterns,
\emph{S\'em. Lothar. Combin.} \textbf{89B} (2023), Article~61, 12~pp.
\href{https://www.mat.univie.ac.at/~slc/wpapers/FPSAC2023/61.pdf}{Available text}.
\bibitem{Pan}
R.~Pan, \emph{Algorithmic Solution to Problem 1
(and linear extensions of general one-level grid-like posets)},
Project P, posted 28 June 2016, pp.~2--3.
\href{https://mathweb.ucsd.edu/~projectp/problems/solutions/OneLevelGridPoset.pdf}{Available text}.
\bibitem{OEIS}
The OEIS Foundation Inc., \emph{The On-Line Encyclopedia of Integer Sequences},
\href{https://oeis.org/A274644}{A274644}, revision 32, 24 October 2026;
recurrence contributed by M.~Wallner, 13 February 2024.
\bibitem{Core}
\emph{BlitvicKammounSteingrimsson14325Core.lean},
\newline revision \texttt{7b7707afe48bfe97846f5f4c44c709e53aec7178}.
\href{https://github.com/the-omega-institute/trureturing/blob/7b7707afe48bfe97846f5f4c44c709e53aec7178/D5/S3/Combinatorics/Permutation/BlitvicKammounSteingrimsson14325Core.lean}{Pinned Lean declarations}.
\bibitem{Transfer}
\emph{BlitvicKammounSteingrimsson14325Transfer.lean}, same revision.
\href{https://github.com/the-omega-institute/trureturing/blob/7b7707afe48bfe97846f5f4c44c709e53aec7178/D5/S3/Combinatorics/Permutation/BlitvicKammounSteingrimsson14325Transfer.lean}{Pinned Lean declarations}.
\end{thebibliography}
\noindent\small Licensed under Creative Commons Attribution 4.0 International
(CC BY 4.0).
\end{document}
